% Generated draft: requires mathematical and source review.
\documentclass[11pt]{article}
\usepackage[margin=25mm]{geometry}
\usepackage{fontspec}
\setmainfont{Latin Modern Roman}
\usepackage{amsmath,amssymb,mathtools,unicode-math}
\setmathfont{Latin Modern Math}
\usepackage{enumitem,needspace,xurl,hyperref}
\setlength{\parindent}{0pt}
\setlength{\parskip}{5pt}
\setlength{\emergencystretch}{3em}
\newcommand{\N}{\mathbb{N}}
\newcommand{\Z}{\mathbb{Z}}
\newcommand{\Q}{\mathbb{Q}}
\newcommand{\R}{\mathbb{R}}
\newcommand{\C}{\mathbb{C}}
\newcommand{\F}{\mathbb{F}}
\newcommand{\A}{\mathbb{A}}
\newcommand{\PP}{\mathbb{P}}
\newcommand{\E}{\mathbb{E}}
\newcommand{\eps}{\varepsilon}
\newcommand{\dd}{\,\mathrm{d}}
\begin{document}
\section*{UnsolvedMath ID: 30006124}
\textbf{Source problem code:} catalogue entry
\subsection{Definitions and mathematical statement}
The supplied target is to prove the weighted fundamental lemma for nonsplit classical groups, with the stated aim of removing the corresponding conditional inputs to the endoscopic classification of their representations. The record does not specify the version of the weighted fundamental lemma, its hypotheses, or its orbital-integral identity. Consequently, it does not determine a complete quantified conjecture. In particular, the assertion that this would make the classification unconditional requires clarification of the other assumptions mentioned in the source.

The following describes the supplied classification context, without asserting an unconditional theorem or supplying the missing lemma. Let \(G\) be a quasi-split symplectic or orthogonal group; the source also discusses quasi-split unitary groups, attributing their treatment to Mok. Its displayed formulas use \(G(\Q_v)\) and \(G(\A)\), where \(v\) is a place of \(\Q\), \(\Q_v\) is the corresponding completion, and \(\A\) is the adele ring of \(\Q\). Whether this base field is intended for the requested lemma is unspecified. Write \(\widehat G\) for the dual group. The source describes \(\Psi(G)\) as a set of formal global parameters constructed using cuspidal automorphic representations of \(\mathrm{GL}_N\), with \(N\) associated with the standard representation of \(\widehat G\). It identifies a subset
\[
\Phi(G)=\{\phi\in\Psi(G):\phi|_{\mathrm{SL}_2}=1\},
\]
whose elements it calls generic, or tempered. The parameter domain and the indicated \(\mathrm{SL}_2\) factor are not defined in the record.

The supplied excerpt attributes to Arthur, under unspecified assumptions, a classification with the following features. There is a decomposition of the discrete automorphic spectrum, written schematically as
\[
L^2_{\mathrm{disc}}(G)=\bigoplus_{\psi\in\Psi_2(G)}L^2_\psi(G).
\]
Here \(L^2_{\mathrm{disc}}(G)\) denotes the source's discrete automorphic spectrum, and \(\Psi_2(G)\) its indexing set of discrete parameters; their precise definitions and the quotient defining the spectrum are not supplied. At each place \(v\), local packets \(\Pi_\phi(G(\Q_v))\) and \(\Pi_\psi(G(\Q_v))\) are described as covering, respectively, the irreducible tempered and irreducible admissible representations. The record does not define the local parameter sets \(\Phi(G_v)\) and \(\Psi(G_v)\), and the damaged transcription does not reliably specify the union or disjoint-union assertions. In particular, the claimed coverage of all irreducible admissible representations by the displayed \(\Psi\)-packets is retained only as a source assertion requiring clarification.

The excerpt further gives a packet-labeling map
\[
\Pi_\psi(G(\Q_v))\longrightarrow\operatorname{Irr}(S_\psi),
\qquad \pi_v\longmapsto\rho_{\pi_v},
\qquad S_\psi=\operatorname{Cent}(\psi,\widehat G).
\]
This is the supplied notation: it does not clarify whether a component group or a quotient of this centralizer is intended. It also gives localization maps \(\operatorname{loc}_v:\Psi(G)\to\Psi(G_v)\) and a schematic multiplicity formula
\[
L^2_\psi(G)=\bigoplus_{\pi}\pi^{\oplus m(\psi,\pi)},
\qquad
m(\psi,\pi)=\operatorname{mult}\!\left(\epsilon_\psi,
\bigotimes_v\left.\rho_{\pi_v}\right|_{S_\psi}\right).
\]
The sum is described as running over irreducible admissible restricted tensor products \(\pi=\bigotimes'_v\pi_v\) of \(G(\A)\) satisfying the source's condition \(\operatorname{loc}_v(\psi)=\psi_{\pi_v}\) at every place. The notation \(\psi_{\pi_v}\), the maps needed for the restrictions, and the tensor-product conventions are unspecified. The source calls \(\epsilon_\psi\) an explicit sign character and says it is trivial for \(\psi\in\Phi(G)\). These formulas record the intended context, not a repaired precise theorem.

The source describes the proof as an induction on the dimension of the standard representation of \(\widehat G\), using comparisons of trace formulas that require local information. It also states an expectation for general groups that cuspidal automorphic representations admitting a Whittaker model have tempered local components and should correspond to \(\Phi(G)\). This is presented as an expectation. Finally, it says that describing all discrete automorphic representations, including those that are not cuspidal, requires the \(\mathrm{SL}_2\) factor even for \(\mathrm{GL}_N\), attributing the conjecture to Arthur and its proof to M\oe glin--Waldspurger. None of this supplies the missing weighted fundamental lemma identity or its precise range of validity.
\subsection{Short English statement}
The requested problem is to establish a weighted fundamental lemma for nonsplit classical groups and thereby remove the corresponding conditional input to endoscopic classification. The supplied record gives conditional classification context, but omits the lemma's identity, hypotheses, and exact scope. A precise theorem statement therefore requires source clarification and human mathematical review. The supplied status is \emph{partially solved}; it is not an independently verified status assessment.
\subsection{Sources}
\begin{itemize}
\item Local UnsolvedMath registry, supplied complete record with permanent ID \(30006124\), exact source problem code \texttt{catalogue entry}, title \emph{Weighted Fundamental Lemma for Classical Groups}. No proposer is named. The supplied proposed year is 2024.
\item Supplied report citation: \emph{Representations of \(p\)-adic Groups} (2025), Oberwolfach Reports, DOI \texttt{10.4171/owr/2024/54}, pp.\ 3103--3178. \url{https://doi.org/10.4171/owr/2024/54}. The record combines this 2025 citation with a 2024 report identifier and proposed year; no reconciliation is supplied.
\item Supplied literature source, ``OWR report'': \url{https://publications.mfo.de/bitstream/handle/mfo/4270/OWR_2024_54.pdf?sequence=1}.
\item Supplied literature source, ``Arthur, Endoscopic Classification'': \url{https://bookstore.ams.org/COLL/61}.
\item Supplied literature source, ``Arthur, endoscopic classification outline'': \url{https://www.claymath.org/library/cw/arthur/pdf/arthur-endoscopic-tifr.pdf}.
\item Literature assessment dated \texttt{2026-08-22}: the supplied curation labels the problem partially solved and states that Arthur's classification establishes the classical-group result under trace-formula and related hypotheses, with unconditional completion still dependent on the stated assumptions. This is supplied dated curation, not independently verified literature assessment. The record itself requests independent verification of time-sensitive status and cited claims.
\item The registry labels its statement \texttt{corrected\_verified} and reports reconstruction and checking against the official Oberwolfach report. These are supplied source-verification claims, not verification performed for this draft. The cited references were not visited in preparing it.
\item The excerpt attributes the symplectic and orthogonal classification to Arthur, the quasi-split unitary treatment to Mok, and the stated discrete-spectrum result for \(\mathrm{GL}_N\) to M\oe glin--Waldspurger following Arthur's conjecture. No separate bibliographic references for Mok or M\oe glin--Waldspurger are supplied.
\end{itemize}
\end{document}
